% Transcription — fonds Alexandre Grothendieck, shelfmark 146, batch 2 (pages 21-40).
% Established from the facsimile archives/batches/146/batch-02.pdf, per the
% skill transcribe-grothendieck.
%
% Pass: Opus 5.5 (claude-opus-5-5), 2026-10-02 — first pass, unchecked against the pages by a human.
%
% Nine of the twenty pages are transcribed: 21, 23, 26, 28, 30, 32, 35, 37, 39.
% The others (22, 24, 25, 27, 29, 31, 33, 34, 36, 38, 40) are blank versos or
% unrelated paper (USTL administrative letters and tables, budget sheets, a
% typed calculus exercise sheet) and are absent.
%
% Content: a presentation by generators and relations of a groupoid attached
% to a configuration of points P^ω_a, Q_a, R_r on a circle (arrows ã_r, b̃_r,
% c̃_r, d̃_r, the loops λ and the "fibre" arrows μ), with the elimination of
% the μ and λ (pages 26-32); then four lemmas on straight-line homotopies
% (1-s)φ_ω ℓ(t) + s ℓ(t) avoiding a complexified line Δ_C, with corollaries
% giving commutative squares in Π_1 V^{**} (pages 35-39).
%
% Dating copied verbatim from src/content/catalogue.ts.

\documentclass[11pt,a4paper]{article}
\input{../preamble/grothendieck}

\folder{146}
\batch{2}
\pages{21}{40}
\foldertitle{Groupoïdes de Teichmüller et les S0 Tg,!ν, S0 TG : notes manuscrites (s.d.).}
\dating{[à partir de 1978]}
\watermark{Édition de démonstration}

\begin{document}

\section{Points du cercle et vecteurs $\widetilde{R}_i$, $\widetilde{Q}^\omega_i$}

\page{21}
\note{la page reprend une figure et des notations qui viennent d'avant le lot ; on ne les a pas consultées.}

\drawing{Un grand cercle, sur lequel sont marqués et nommés des points : en haut $\widetilde{R}^\omega_k$, $\widetilde{R}_k$, $\widetilde{R}^{\omega'}_k$, puis en tournant $\widetilde{Q}^{\omega'}_j$, $-\widetilde{R}_i$ et ses deux voisins, $\widetilde{Q}^\omega_k$, $\widetilde{R}_j$ et ses deux voisins, $\widetilde{Q}^{\omega'}_i$, en bas $-\widetilde{R}^{\omega'}_k$, $-\widetilde{R}_k$, $-\widetilde{R}^\omega_k$, puis $\widetilde{Q}^\omega_j$, $\widetilde{R}_i$ et ses voisins, $\widetilde{Q}^{\omega'}_k$, $-\widetilde{R}_j$ et ses voisins, $\widetilde{Q}^\omega_i$. Chacun des six points $\pm\widetilde{R}$ est entouré d'un petit demi-cercle ; deux triangles inscrits forment une étoile à six branches, et quelques traits au crayon partent du centre et vers l'extérieur.}

\marginal{où $\omega_{ij}$ est l'orientation telle que $\omega_{ij}(j) = i$}

$\{i,j,k\} = {}$\uncertain{$J$} :
\[
\widetilde{R}_i = (e_j - e_k)\wedge\omega_{jk}, \qquad
\widetilde{Q}^\omega_i = \widetilde{R}_{\omega^2 i} - \widetilde{R}_{\omega i},
\]
donc
\[
i_\omega \widetilde{Q}^\omega_i = \widetilde{R}_i, \qquad
i_\omega \widetilde{R}_i = -\widetilde{Q}^\omega_i = \widetilde{Q}^{\omega'}_i ,
\]
\[
\widetilde{Q}^\omega_i = (e_j + e_k - 2e_i)\wedge\omega .
\]
\note{l'exposant de $\widetilde{R}_i$ dans la première formule est surchargé et illisible ; l'indice de $\widetilde{R}_{\omega^2 i}$ est corrigé.}

\page{23}
\drawing{Page de figures, sans texte suivi. En haut, un cercle avec une flèche courbe notée $3$ (rotation d'ordre 3), deux triangles inscrits, des points marqués de part et d'autre des sommets et étiquetés de petits symboles $R$, $Q$ avec exposants $2$, $3$, $\bar 3$ (illisibles dans le détail) ; des tirets vont d'un sommet vers le centre. En bas à gauche, une sphère avec équateur et méridien, les points $Q_+$ (en haut), $Q_-$ (en bas), $P_+$, $Q_0$, $P'$ sur l'équateur, avec de petites boucles autour de certains points. Au milieu, une petite ellipse et un petit cercle partagé en trois. En bas à droite, un tétraèdre, une flèche, puis une sphère traversée par une ellipse avec des flèches tangentes.}

\section{Générateurs et relations}

\page{26}
\note{$r$ parcourt $R$, et $r = (s,a,\omega)$ ou $(s,a,\omega)$ selon l'ordre adopté ; $r' = \sigma_0 r = (s,a')$.}

Flèches \uncertain{« obliques »} (en face, les objets $P^\omega \leftarrow \widetilde{P}^\omega_a$, $Q_a \leftarrow \widetilde{Q}_a$, $R_r \leftarrow \widetilde{R}_r$, indexés par $(a,\omega) \in A\times\underline{\omega} \simeq R$, $a\in A$, $r\in R$) :
\[
\begin{array}{lll}
a_r \leftrightarrow \tilde a_r : \widetilde{R}_r \to \widetilde{R}_{r'} & & r \in R\\
b_r \leftrightarrow \tilde b_r : \widetilde{Q}_a \to \widetilde{R}_r & & r\in R\\
c_a \leftrightarrow \tilde c^\omega_a : \widetilde{P}^\omega_a \to \widetilde{Q}_a \ \text{ou}\ \tilde c_r & & (a,\omega)\in A\times\underline{\omega}\simeq R
\end{array}
\]

Flèches « fibres » \uncertain{verticales} :
\[
\begin{array}{lll}
1 \leftrightarrow \lambda_{\widetilde{P}^\omega_a}\ \text{lacet en}\ \widetilde{P}^\omega_a & & (a,\omega)\in A\times\underline{\omega}\simeq R\\
1 \leftrightarrow \lambda_{\widetilde{Q}_a}\ \text{lacet en}\ \widetilde{Q}_a & & a \in A\\
1 \leftrightarrow \lambda_{\widetilde{R}_r}\ \text{lacet en}\ \widetilde{R}_r & & r\in R\\
\mu_{\widetilde{Q}_a} : \widetilde{Q}_a \to \widetilde{Q}_a^\vee & & a\in A\\
\mu_{\widetilde{R}_r} : \widetilde{R}_r \to \widetilde{R}_r^\vee & & r\in R\\
1 \leftrightarrow \tilde d^\omega_r : \widetilde{P}^\omega_{a_1} \to \widetilde{P}^\omega_{\bullet}\ \text{ou}\ \tilde d^\omega_s & & r = (s,\omega)\in S\times\underline{\omega}\simeq R
\end{array}
\]
\note{une flèche courbe porte « renversons l'ordre » entre $\widetilde{P}^\omega_a$ et $\widetilde{R}_r$, une autre entre $\mu_{\widetilde{Q}_a}$ et $\mu_{\widetilde{R}_r}$ ; à gauche, une colonne de symboles biffés ($\mu$, $\lambda$) que l'auteur a remplacés par $1$. Les indices de droite de la colonne des $\mu$ sont en partie biffés : $(a,\ldots)\in A$.}

« relations \ill{} » :
\[
(1)\quad \tilde a_r \tilde b_r \tilde c_r \tilde d_r = \tilde b_{r_1} \tilde c_{r_2}
\qquad r\in R,\ r' = \sigma_0 r,\ r_1 = \sigma_1\sigma_0(r') = \omega^{-1} r'
\]
« relations fibres » :
\[
(2)\quad \tilde d_{\omega^m r}\cdots \tilde d_{\omega r}\tilde d_r = \lambda_{\widetilde{P}^\omega_{a_1}}
\qquad r\in R,\ r = (s,a,\omega),\ a' = \omega^{-1}(a)
\]
\[
(6)\quad \mu_{\widetilde{Q}_a^\vee}\,\mu_{\widetilde{Q}_a} = \lambda_{\widetilde{Q}_a}
\qquad (7)\quad \mu_{\widetilde{R}_r^\vee}\,\mu_{\widetilde{R}_r} = \lambda_{\widetilde{R}_r}
\qquad r\in R
\]
\note{les numéros des relations fibres sont surchargés : $(2)$ corrigé en $(5)$, $(5)$ en $(6)$, $(6)$ en $(7)$ ; on donne la dernière lecture là où elle se laisse lire.}

« relations de transport » ou « de commutation » :
\[
\begin{array}{lll}
(8) & \tilde a_r(\lambda_{\widetilde{R}_r}) = \lambda_{\widetilde{R}_{r'}} & r\in R,\ r' = \sigma_0(r)\\
(9) & \tilde b_r(\lambda_{\widetilde{Q}_a}) = \lambda_{\widetilde{R}_r} & r\in R,\ r = (a,s,\omega)\\
(10) & \tilde c_r(\lambda_{\widetilde{P}^\omega_a}) = \lambda_{\widetilde{Q}_a} & (a,\omega)\in A\times\underline{\omega}\simeq R
\end{array}
\]

\page{28}
\struck{$\mu_{\widetilde{Q}^{\omega}_a{}^\vee}\,\mu_{\widetilde{Q}^\omega_a} = \lambda_{\widetilde{Q}^\omega_a} = \tilde e^\omega_a(\lambda_{\widetilde{P}^\omega_a})$} \quad $\lambda_{\widetilde{Q}^\omega_a}$ — \marginal{(2), entouré}

\struck{$\widetilde{A}_{\omega^n r^*}\cdots\widetilde{A}_{\omega r^*}\widetilde{A}_{r^*}\,\tilde c_r = \tilde c_{r^\vee}\,\tilde d_{\omega^n r^*}\cdots\tilde d_{\omega r^*}\tilde d_{r^*}$}
\quad conséquence de (1)
\note{encadré puis barré de traits obliques ; une accolade sous le produit des $\widetilde{A}$ porte « déf $= \mu_{\widetilde{Q}^\omega_a}$ flèche fibre ».}

\[
(3)\quad \mu_{\widetilde{R}_r}\circ\tilde b_r = \tilde b_{r^\vee}\circ\mu_{\widetilde{Q}^\omega_a}
\]
\struck{i.e. l'autre}
\[
(4)\quad \mu_{\widetilde{R}_{r'}}\circ\tilde a_r = \tilde a_{r^\vee}\,\mu_{\widetilde{R}_r}
\]
\[
(2)\quad \mu_{\widetilde{Q}_a}\,\tilde c_r = \tilde c_{r^\vee}\,\tilde d_{\omega^n r^*}\cdots\tilde d_{r^*}
\ \Longrightarrow\ \ill{}
\]
\[
\overset{\text{via (1)}}{\Longrightarrow}\quad
\mu_{\widetilde{Q}_0} = \widetilde{A}_{\omega^n r^*}\cdots\widetilde{A}_{\omega r^*}\widetilde{A}_{r^*},
\qquad \widetilde{A}_r = \tilde b_r^{-1}\tilde a_r^{-1}\tilde b_{r_1}
\]
d'où $\mu_{\widetilde{R}_r}$ via (3) :
\[
\mu_{\widetilde{R}_r} = \widetilde{B}_{\omega^n r}\cdots\widetilde{B}_{\omega r}\widetilde{B}_r
\quad\text{où}\quad
\widetilde{B}_r \overset{\text{déf}}{=} \tilde a_{r''}^{-1}\,\tilde b_{r'}\,\tilde b_r^{-1}
\]
et (4) devient une relation entre les $\tilde a_r, \tilde b_r, \tilde c_r$, à laquelle il faut ajouter l'expression \uncertain{de} $\tilde c_r$ \ill{} \ill{} et $\mu_{\widetilde{Q}_0}$ \ill{}.
\marginal{2 groupes de relations \ill{} entre $\tilde a_r$, $\tilde b_r$ \ill{} ($\mu$ éliminés \ill{}) ; \ill{} $\in A$, $r\in R$ ($\mu_{\widetilde{Q}^\omega_0} = \mu_{\widetilde{Q}_0}$ \ill{}) \ill{} la relation (5) \ill{}}

Mais on voit que \uncertain{ces deux relations} \ill{} \uncertain{équivalentes} et $\pm$ triviales.

On peut considérer (5), (6), (7) comme \emph{définitions} des $\lambda_{\widetilde{P}^\omega_\bullet}$, des $\lambda_{\widetilde{Q}_0}$ et $\lambda_{\widetilde{R}_r}$ en termes des $\tilde a_r, \tilde b_r, \tilde c_r$) dans \uncertain{les} équations $(8, 9, 10)$ \uncertain{deviennent} des équations en $\tilde a_r, \tilde b_r, \tilde c_r$.

\page{30}
\[
(2')\quad \widetilde{A}_{\omega^m r^{**}}\cdots\widetilde{A}_{r^{**}} = \widetilde{A}_{\omega^m r_1^*}\cdots\widetilde{A}_{r_1^*}
\qquad \forall a\in A
\]
\marginal{(3), entouré}

les relations de transport $(10)$ \struck{\ill{}} \uncertain{identique}, $(9)$ \ill{}, $(8)$ \ill{}

Il reste
\[
\begin{array}{ll}
(1) & \tilde a_r\tilde b_r\tilde c_r\tilde d_r = \tilde b_{r_1}\tilde c_{r_2}\\[2pt]
 & \quad r\in R,\ r_1 = \sigma_0(r),\ r_2 = \sigma_1\sigma_0(r) = \omega^{-1}(r)\\[4pt]
(2') & \widetilde{A}_{\omega^m r''}\widetilde{A}_{\omega^{m-1} r''}\cdots\widetilde{A}_{r''} = \widetilde{A}_{\omega^m r_1}\cdots\widetilde{A}_{r_1}\\[2pt]
 & \quad r\in R,\ r_1 = \sigma_0 r,\ r'' = \sigma_0\sigma_1(r) = \omega(r)
\end{array}
\]
où on pose $\widetilde{A}_r = \tilde b_r^{-1}\tilde a_r^{-1}\tilde b_{r_1}$ ; le membre de gauche de (2') est, par définition, $\mu_{\widetilde{Q}_a}$.
\note{dans la marge gauche, un passage surchargé et biffé, où se lisent « (2') », « 6 relations » et $\widetilde{A}_r$ ; illisible pour le reste.}

introduction des $\mu$ :
\[
\begin{aligned}
\mu_{\widetilde{Q}_0} &\overset{\text{déf}}{=} \widetilde{A}_{\omega^m r''}\cdots\widetilde{A}_{r''} = \cdots\\
\mu_{\widetilde{R}_r} &\overset{\text{déf}}{=} \widetilde{B}_{\omega^m r}\cdots\widetilde{B}_r
\qquad\text{où}\quad \widetilde{B}_r \overset{\text{déf}}{=} \tilde a_{r''}^{-1}\tilde b_{r'}\tilde b_r^{-1}\\
\mu_{\widetilde{P}^\omega_a} &\overset{\text{déf}}{=} \tilde d_{\omega^m r''}\cdots\tilde d_{\omega r''}\tilde d_{r''}
\end{aligned}
\]
$\Longrightarrow$ relations (3) (4) (5) de « commutation » de ces flèches avec $\tilde a$, $\tilde b$, $\tilde c$.

introduction des $\lambda$ :
\[
\begin{aligned}
\lambda_{\widetilde{Q}_0} &\overset{\text{déf}}{=} \mu_{\widetilde{Q}_a^\vee}\,\mu_{\widetilde{Q}_a} = \widetilde{A}_{\omega^m r''}\cdots\widetilde{A}_{r''} = \widetilde{A}_{\omega^m r_1}\cdots\widetilde{A}_{r_1}\\
\lambda_{\widetilde{R}_r} &\overset{\text{déf}}{=} \mu_{\widetilde{R}_r^\vee}\,\mu_{\widetilde{R}_r} = \widetilde{B}_{\omega^m r}\cdots\widetilde{B}_r\\
\lambda_{\widetilde{P}^\omega_a} &\overset{\text{déf}}{=} \mu_{\widetilde{P}^\omega_{a^\vee}}\,\mu_{\widetilde{P}^\omega_a} = \tilde d_{\omega^m r''}\cdots\tilde d_{r''}
\end{aligned}
\]
\struck{relations} de \uncertain{transport} (8), (9), (10) \ill{} \quad Commutation des flèches $\tilde a_r, \tilde b_r, \tilde c_r$ avec \ill{} $\lambda$)
\note{l'exposant $m$ ou $n$ des produits n'est pas tenu fixe d'une ligne à l'autre ; on le transcrit tel qu'il se lit.}

\page{32}
\[
(2'')\quad \tilde c^\omega_a(\lambda_{\widetilde{P}^\omega_a}) = \tilde c^\omega_a(\lambda_{\widetilde{P}^{\omega'}_a})
\]
\marginal{(4)}
en présence de (1), c'est une conséquence de (2'), qui équivaut à la relation déduite de (2') en \uncertain{la portant au carré}.

\begin{tikzcd}[column sep=large]
Q \arrow[r, bend left=40, "\mu'"] \arrow[r, bend left=15, "\mu"'] & Q^\vee \arrow[l, bend left=15, "\check\mu"'] \arrow[l, bend left=40, "\check\mu'"]
\end{tikzcd}

\note{le sens des quatre flèches est lu sur la figure : $\mu, \mu'$ vont de $Q$ à $Q^\vee$, $\check\mu, \check\mu'$ en reviennent ; les têtes sont peu nettes.}

\[
\underbrace{\check\mu\,\mu}_{\lambda} = \underbrace{\check\mu'\,\mu'}_{\lambda'}
\qquad
\check\mu'^{-1}\check\mu = \mu'\mu^{-1} \quad (\mu'^{-1}\mu)^{-1}
\]
\[
\mu' = 1,\ \check\mu' = \check\mu \quad\Longleftarrow\quad \check\mu = \check\mu'\mu'
\]

\[
\overbrace{\widetilde{A}_{\omega^n r''}\widetilde{A}_{\omega^{n-1} r''}\cdots\widetilde{A}_{r''}}^{\mu_{\widetilde{Q}^\omega_a}}
= \overbrace{\widetilde{A}_{\omega^n r_1}\widetilde{A}_{\omega^{n-1} r_1}\cdots\widetilde{A}_{r_1}}^{\mu_{\widetilde{Q}^{\omega'}_{\bullet}}}
\]
équivaut à : ceci, \struck{\ill{}} \struck{$\widetilde{B}_{\omega^m r}\widetilde{B}_{\omega r}\cdots\widetilde{B}_r = \mu_r$}

\drawing{Une chaîne de points sur une courbe fermée, chacun muni de petites boucles fléchées (les lacets $\lambda$), reliés par des flèches à double sens. Au premier plan, en trait gras : $\widetilde{Q}_a \xrightarrow{\tilde b_r} \widetilde{R}_r$, puis un arc $\tilde a_r$ de $\widetilde{R}_r$ vers $\widetilde{R}_{r_1}$, sous lequel un arc $\tilde a_{r_1}$ ; au-dessus, un autre point avec l'arc $\tilde a_{r_1^\vee}$. La région entre $\widetilde{R}_r$, $\widetilde{Q}_a$ et un sommet commun est hachurée ; des traits fins portent $r$, $r_1$, $r_2$. Plus bas à droite, un petit triangle isolé.}

\section{Lemmes sur les homotopies linéaires}

\page{35}
\textbf{Lemme 1.} Soit $\overrightarrow{[x,y]}$ un \uncertain{segment orienté} \struck{\ill{} tel que la droite \ill{}} \struck{$\Delta$ \ill{} droite \ill{} qui \ill{}} $\Delta_0$ \ill{}, \uncertain{orienté}, \uncertain{ne contienne pas l'origine}, \struck{\ill{} $x, y$, on suppose que} \ill{} \uncertain{orientation} \ill{} de $\Delta$ \struck{\ill{}} (tel que l'orientation de $\Delta$ \uncertain{induite} \uncertain{par la} \ill{} \uncertain{aille} de $x$ vers $y$).
\note{les premières lignes du lemme sont plusieurs fois reprises et biffées ; on n'en donne que ce qui se lit.}

Soit $\Delta$ une droite \uncertain{horizontale} \struck{\ill{}} telle que $\Delta\cap\{x,y\} = \emptyset$. Considérons
\[
f_{x,y}(s,t) = (1-s)\,\varphi_\omega\,\ell_{y,x}(t) + s\,\ell_{y,x}(t),
\]
qui, pour $t$ fixé, parcourt le segment $\varphi_\omega\ell_{y,x}(t) : \ell_{y,x}(t)$. Alors $f_{x,y}([0,1]\times[0,1])$ ne rencontre pas $\Delta_{\mathbb C}$.

\marginal{$\ell_{y,x}(t) = \frac{x+y}{2} + e(t)\frac{x-y}{2}$, \quad $e(t) = \exp i\pi t$.}
\marginal{$L_{y,x}(t) = \uncertain{(1-t)x + ty}$}
\note{en marge gauche, une note oblique de plusieurs lignes, en partie biffée, commençant par « ici » : illisible.}

\textbf{Corollaire.} Si $x, y\in V^{**}$ \struck{\ill{}} (cf. \ill{}) \uncertain{i.e.} $x, y$ correspondent à \ill{} de $\Sigma^*$) dans $\Pi_1 V^{**}$, \ill{} :

\begin{tikzcd}[column sep=huge, row sep=large]
x \arrow[r, "\ell_{y,x}"] \arrow[d, leftrightarrow, "L_{x,\varphi_\omega(x)}"'] & y \arrow[d, leftrightarrow, "L_{y,\varphi_\omega(y)}"] \\
\varphi_\omega(x) \arrow[r] & \varphi_\omega(y)
\end{tikzcd}

\[
\varphi_\omega(\ell_{y,x}) = \ell_{\varphi_\omega(y),\varphi_\omega(x)}
\]
\note{les flèches verticales ont une tête à chaque bout sur la page.}

\page{37}
\textbf{Lemme 2} (trivial). \struck{\ill{}} Posons, pour $x, y\in V$ et un chemin $\ell : x\to y$ dans $V$,
\[
g(s,t) = \text{\struck{$1$}}\ (1-s)\,\varphi_\omega\,\ell(t) + s\,\ell(t).
\]
\note{le facteur $s$ du second terme ne se distingue pas nettement ; on le lit par analogie avec le lemme 1.}
Soit $\Delta$ une droite de $V$ telle que \struck{$\ell[0,1]$} $|\ell|\cap\Delta = \emptyset$. Alors $g([0,1]\times[0,1])\cap\Delta_{\mathbb C} = \emptyset$.
En effet, $\Re\, g([0,1]\times[0,1]) = |\ell|$, et $\Re\,\Delta_{\mathbb C} = \Delta$.

\textbf{Corollaire.} \struck{\ill{}} Si $\ell$ est un chemin dans $V^{**}$, \ill{} \uncertain{commutatif} dans \ill{} :

\begin{tikzcd}[column sep=huge, row sep=large]
\varphi_\omega(x) \arrow[r, "\varphi_\omega(\ell)"] \arrow[d, "L_{x,\varphi_\omega(x)}"'] & \varphi_\omega(y) \arrow[d, "L_{y,\varphi_\omega(y)}"] \\
x \arrow[r, "\ell"] & y
\end{tikzcd}

\textbf{Lemme 3} ($\pm$ trivial). $|\ell_{y,x}|\cap\Delta_{\mathbb C} = \emptyset$ pour \uncertain{toute} droite $\Delta$ qui ne rencontre pas $\{x,y\}$.
\note{sous « Lemme 3 », une ligne courte entre parenthèses, illisible.}

Car on a $\Re\,\ell_{y,x} = [x,y]$, \ill{} \ill{} $|\ell_{y,x}|\cap\Delta_{\mathbb C}\neq\emptyset$ \uncertain{pour} \ill{} ($\Delta\cap\{x,y\} = \emptyset$) \uncertain{donc} $\Delta\cap[x,y]\neq\emptyset$, \ill{}, $\Delta\cap[x,y] = \{a\}$, $a\neq x,y$, \ill{}, $a = \Re\,\ell_{x,y}(t)$, \uncertain{avec} $t\neq 0,1$ ; \uncertain{alors} $\operatorname{Im}\ell_{x,y}(t) = (\sin t)\,\frac{y-x}{2}$ dans \ill{} $y - x\in$ \ill{} \ill{} $\Delta$ \ill{} $x, y\in\Delta$ \ill{}.
\note{le $\sin t$ (sans le facteur $\pi$ de $e(t) = \exp i\pi t$, page 35) est lu tel qu'écrit.}

\page{39}
\textbf{Lemme 4.} Posons
\[
\ell_x(t) = \ell_{-x,x}(t) = e(t)\,x : x\to -x,
\]
et considérons
\[
\begin{aligned}
h(s,t) &= (1-s)\ \text{\struck{$\cdots$}}\ \ell_{\varphi(x)}(t) + s\,\ell_x(t)\\
&= e(t)\bigl((1-s)\varphi(x) + s x\bigr).
\end{aligned}
\]
Alors pour toute droite $\Delta$ de $V$ telle que $x\notin\Delta$, on a
\[
h([0,1]\times[0,1])\cap\Delta_{\mathbb C} = \emptyset.
\]
En effet, $h(s,t)\in\Delta_{\mathbb C}$ $\Longleftrightarrow$ $(1-s)\varphi(x) + s x\in\Delta_{\mathbb C}$ $\Longrightarrow$ $(1-s)x + s x\in\Delta$ i.e. $x\in\Delta$ \uncertain{pour} \ill{} $\Re$.
\note{dans la définition de $h$, un premier essai est biffé et illisible ; à la dernière ligne, l'implication passe à la partie réelle ; le mot qui précède $\Re$ est illisible.}

\textbf{Corollaire.} \uncertain{Le} \ill{} \uncertain{commutatif} dans $\Pi_1 V^{**}$, si $x\in V^{**}$ :

\begin{tikzcd}[column sep=huge, row sep=large]
\varphi_\omega(x) \arrow[r, "\ell_{\varphi_\omega(x)}"] \arrow[d, "L_{x,\varphi_\omega(x)}"'] & -\varphi_\omega(x) \arrow[d, "L_{-x,\varphi_\omega(x)}"] \\
x \arrow[r, "\ell_x"] & -x
\end{tikzcd}

\note{l'étiquette de droite se lit $L_{-x,\varphi_\omega(x)}$, et non $L_{-x,-\varphi_\omega(x)}$ ; on la laisse telle. La suite des lemmes déborde peut-être au-delà du lot.}

\end{document}
